\section{模型评价与结论}

\subsection{模型优点}

第一，本文将“连续重复条目”和“有效到达”明确区分，任务展开、路线构造和验证器使用同一口径，避免把 A-A-A 错计为三次巡检。第二，三个问题沿同一任务集合和基地闭环约束递进：问题一确定机队和最大工作时间，问题二只做固定机队内的均衡优化，问题三再加入时间依赖禁飞区，模型继承关系清晰。第三，问题三同时检查几何和时间，不用固定比例增加距离代替真实绕行；等待、绕行、服务和基地状态均进入独立验证。第四，结果保存为 JSON 单一真源，图表和表格均由真实运行结果生成，便于复现。

\subsection{模型不足}

第一，主要求解器是固定随机种子多起点启发式，本轮没有精确 MIP 最优性或不可行性证书，因此不能把首个可行机队称为全局最小。第二，圆形禁飞区的复杂交叠场景需要可见图和时间依赖边搜索，当前正式 Case 采用可行优先的局部搜索，仍可能存在未发现的更优路线。第三，模型忽略风场、起降耗时、电池续航和无人机之间的空域冲突，适用于题面给定的静态速度和任务约束，推广到真实低空运行时需要补充这些因素。

\subsection{主要结论}

图~\ref{fig:casecompare} 从机队、最大工作时间和均衡指标三个层面汇总了四个 Case 在三个问题中的真实结果，便于观察模型递进后的变化。

\begin{figure}[H]
\centering
\includegraphics[width=0.92\textwidth]{case_comparison.pdf}
\caption{四个 Case 在三个问题中的真实结果对比}
\label{fig:casecompare}
\end{figure}

问题一的最好已验证机队数为 Case1--Case4 的 5、2、6、5 架；对应最大工作时间为 8.7129、7.9583、8.6125、8.1378 h，均通过 $9\ \mathrm{h}$ 校验。问题二保持相同机队，最大工作时间为 8.6139、7.8011、8.5980、8.1378 h，极差为 0.0896、0.0027、0.0976、0.0644 h，说明跨路线任务调整有效改善了均衡性。

问题三在动态禁飞区下的最小已验证可行机队数为 5、3、9、7 架。Case2 需要 1.503 h 安全等待，Case3 需要 0.454 km 绕行；固定原机队的 Case2--Case4 对照方案分别达到 10.1388、10.4125、10.3717 h，因此必须增加无人机才能恢复 $9\ \mathrm{h}$ 可行性。四个正式 q3 结果均通过任务覆盖、有效到达、基地闭环、时间传播和禁飞约束的独立复核。

总体而言，本模型给出了一条从“任务展开”到“机队确定”、从“负载均衡”到“动态禁飞时空重规划”的统一解决路径。其结果适合作为题面参数下的最好已验证调度方案；若需要宣称全局最优或对参数扰动做定量敏感性结论，还应补充精确求解和专门的灵敏度实验。
